\chapter{Discusión}
\epigrafeen{The first principle is that you must not fool yourself---and
you are the easiest person to fool.}{Richard P. Feynman}{Cargo Cult
Science, 1974}

\pendiente{redactar el resto de la discusión}

\section{Qué mide un control negativo}
\label{sec:disc_controles}

Los experimentos de control negativo no aportan información sobre el efecto del campo magnético. Aportan información sobre el aparato que se emplea para medirlo, y esa distinción merece explicitarse, porque determina el peso que puede asignarse a los resultados de los tratamientos.

Feynman \cite{feynman_1974} relata un experimento que ilustra el punto con precisión. En 1937, Young dispuso un corredor con puertas a ambos lados y se propuso entrenar ratas para que entraran siempre por la tercera puerta a partir del punto de partida. Las ratas, en cambio, se dirigían directamente a la puerta donde había estado el alimento la vez anterior. Dado que el corredor era uniforme, algo debía distinguir a esa puerta de las demás. Young pintó todas las puertas replicando exactamente sus texturas; las ratas seguían acertando. Supuso entonces que olfateaban el alimento y alteró los olores con reactivos después de cada corrida; las ratas seguían acertando. Consideró que podían orientarse por las luces y la disposición del laboratorio, y cubrió el corredor; las ratas seguían acertando. Finalmente descubrió que se orientaban \textbf{por el sonido que producía el piso bajo sus patas}, y solo consiguió suprimir esa pista apoyando el corredor sobre arena. Recién entonces las ratas debieron aprender la consigna. Si relajaba cualquiera de esas condiciones, volvían a acertar.

La observación de Feynman sobre este trabajo es la que interesa aquí: Young no descubrió nada acerca de las ratas, pero descubrió todo lo que hay que controlar para poder descubrir algo acerca de las ratas. El autor señala además que los experimentos posteriores ignoraron sistemáticamente ese resultado y continuaron corriendo ratas del mismo modo de siempre, sin citarlo, precisamente porque no constituía un hallazgo sobre el objeto de estudio.

Los controles negativos de este trabajo cumplen esa función. El análisis conjunto de los tres experimentos (Sección~\ref{subsec:ctrl_negativo}) reveló un gradiente vertical significativo entre los pisos extremos del sistema de germinación, y una anomalía localizada en la celda C9, ambos reproducibles con las tres métricas de crecimiento evaluadas. Ninguno de los dos es un resultado sobre campos magnéticos. Los dos son resultados sobre el montaje: son el equivalente del sonido del piso en el corredor de Young. Sin haberlos buscado deliberadamente, una diferencia de esa magnitud podría haberse atribuido al tratamiento, dado que su tamaño es comparable al de los efectos que la literatura reporta para el magnetopriming.

Conviene precisar qué protección ofrece este conocimiento y cuál no. Detectar el gradiente no lo elimina: lo que impide que se confunda con el efecto del tratamiento es la asignación aleatoria de las muestras a las posiciones de germinación (Sección~\ref{subsec:asignador}), que distribuye esa variabilidad entre condiciones en lugar de dejarla asociada a alguna de ellas. La caracterización y la aleatorización operan, por lo tanto, de manera complementaria: la primera establece que el sesgo existe y cuánto vale; la segunda impide que se alinee con la variable de interés.

Corresponde, por último, aplicar el mismo criterio al propio diseño. La torre de exposición no cuenta con un control equivalente: el nivel de apilamiento determina simultáneamente la dosis magnética y la altura, de modo que ambos factores permanecen confundidos durante la hora de exposición. En los términos de la analogía anterior, se caracterizó el piso del corredor de germinación pero no el de la torre. Un experimento consistente en apilar los cajones sin imán, replicando el resto del protocolo, es la verificación que falta, y es exactamente el tipo de experimento que Feynman señala que la literatura tiende a no realizar y a no citar.

\endinput
